\subsection{Financial Time-Series Synthesis}
\label{sec:synthesis}

\subsection{Problem Formulation and Characteristics}

Financial Time-Series Synthesis (FTSS) refers to the systematic, algorithmic generation of statistically credible and temporally coherent synthetic sequences. The overarching objective is to enlarge or enrich the empirical sample space upon which downstream quantitative models and trading agents are trained. Unlike extrapolation (which predicts the continuation of a partial series) or imputation (which mathematically restores broken or missing intervals), FTSS is designed to create entirely new, unobserved data manifolds from scratch. Formally, the goal is to learn a generative model $p_\theta(\mathbf{X})$ that closely approximates the true multidimensional data-generating distribution $p_{data}(\mathbf{X})$, enabling unconditional or condition-guided sampling of synthetic market trajectories $\hat{\mathbf{X}}$.

The primary functions of FTSS include restoring population heterogeneity, rebalancing extremely heavy-tailed historical episodes to correct dataset skewness, and providing controlled, counterfactual market regimes for rigorous stress testing. As shown in Figure \ref{fig:ftss_comparison}, methodologically, the literature can be categorized across two distinct scales: \textit{micro-level synthesis}, which focuses on generating univariate or low-dimensional asset prices and return sequences, and \textit{macro-level market simulation}, which attempts to reproduce complex, asynchronous market microstructures such as order flows and full Limit Order Books (LOB).




\subsection{Micro-Level Synthesis: Asset Returns and Pricing Dynamics}

Micro-level synthesis targets the generation of asset price trajectories and return sequences, serving primarily as a sophisticated data augmentation mechanism. Early applications of standard Generative Adversarial Networks (GANs) in this domain frequently suffered from severe training instability and mode collapse, largely due to the unique "stylized facts" of financial data, such as pronounced volatility clustering and non-Gaussian heavy tails. 

To overcome these barriers, the field rapidly transitioned to Wasserstein GANs (WGAN) \citep{naritomi2020data}, which replace the structurally unstable Jensen-Shannon divergence with the Wasserstein-1 (Earth Mover's) distance. The optimization objective focuses on minimizing:
\begin{equation}
W_1(\mathbb{P}_r, \mathbb{P}_g) = \inf_{\gamma \in \Pi(\mathbb{P}_r, \mathbb{P}_g)} \mathbb{E}_{(\mathbf{x}, \tilde{\mathbf{x}}) \sim \gamma}[\|\mathbf{x} - \tilde{\mathbf{x}}\|]
\end{equation}
where $\Pi(\mathbb{P}_r, \mathbb{P}_g)$ denotes the set of all joint distributions mapping the real distribution $\mathbb{P}_r$ to the generated distribution $\mathbb{P}_g$. The WGAN family was extended for finance in several methodologically distinct directions, which we now unpack critically.

QuantGAN \citep{wiese2020} established the WGAN+TCN template for multi-step unconditional price-path generation. Its verification protocol relies on reproducing eight classical stylized facts (leverage, volatility clustering, heavy tails, absence of autocorrelation in returns, long-memory in absolute returns, etc.) on synthetic S\&P~500 paths; subsequent work \citep{wiese2020} reports that QuantGAN-generated trajectories pass all eight checks while the baseline convolutional GAN fails on at least three. QuantGAN is, however, restricted to \emph{univariate} return series---a limitation that every constraint-aware variant below explicitly tries to relax.

RSQGAN \citep{RSQGAN} injects \emph{regime awareness} into the generator by conditioning on a latent regime variable (typically inferred via a hidden Markov model over historical returns). The reported value proposition is distributional alignment \emph{across} regimes rather than within a single regime: paths sampled under the bull regime should resemble bull-period historical paths, and likewise for bear and transition periods. The conditioning mechanism is comparatively cheap (a single discrete regime index appended to the noise input) and is the only one of the four that does not require a custom loss function.

DAT-CGAN \citep{DAT-CGAN} reformulates the GAN objective around \emph{decision-related} quantities rather than marginal distributions. Concretely, the discriminator is trained against a multi-Wasserstein loss evaluated on portfolio-level statistics (e.g.\ cumulative PnL, Sharpe, drawdown) computed on rolling windows of the generated path. The authors demonstrate on a multi-time-step portfolio-choice task that DAT-CGAN improves generative quality on both the underlying data and on decision-related quantities relative to strong GAN baselines. The cost is a loss function that is harder to optimize: training stability degrades when the decision-statistic set grows.

Jinkou \citep{Jinkou} imposes \emph{state-variable constraints} that keep generated trajectories consistent with economically meaningful drivers (e.g.\ macro factors, implied volatility surfaces). Unlike the other three, Jinkou is positioned for scenario generation rather than return-path generation, which makes it a natural bridge from micro-level synthesis to macro-level market simulation.

TAIL-GAN \citep{cont2025tailgan} stands apart from the other three because its training criterion is not a divergence over the data distribution at all but a \emph{risk-functional score} based on the joint elicitability of Value-at-Risk (VaR) and Expected Shortfall (ES). The discriminator is optimized against the Acerbi--Szekely scoring function, which is the standard backtesting score for VaR/ES; the generator therefore receives gradient signal dominated by tail outcomes rather than by typical samples. This matters because, as \citep{cont2025tailgan} argue explicitly, Wasserstein- or cross-entropy-based losses are sample-dominated and structurally fail to capture tail events---which is exactly the regime regulators care about. Their empirical evaluation compares Tail-GAN against a Historical Simulation baseline and a plain WGAN on a five-name universe (AAPL, AMZN, GOOG, JPM, QQQ) for the period 2019-11 to 2019-12. On out-of-sample VaR/ES estimation for static buy-and-hold, mean-reversion, and trend-following strategies, the reported OOS relative errors are: Historical Simulation $\sim$10.4\%, WGAN $\sim$26.9\%, Tail-GAN $\sim$10.1\%---i.e.\ Tail-GAN closes most of the gap to the oracle, while a vanilla WGAN does roughly 2.5$\times$ worse. On the more demanding score-based test of distributional fidelity under dynamic strategies, Tail-GAN achieves a 21.3\% rejection rate versus WGAN's 11.4\% and historical simulation's 0\%---the higher rejection rate being preferable, since it indicates the synthetic distribution is distinguishable from the test set in the way real markets are. Two limitations temper the result: (i) the headline numbers are on a one-month training window with five names, so scalability to a full S\&P~500 universe required eigenportfolio-based dimensionality reduction, with Tail-GAN(Eig) at 25.6\% OOS error versus Historical Simulation at 25.9\% (a much smaller gap); (ii) Tail-GAN does \emph{not} target stylized-fact reproduction---its evaluation is risk-functional only, so it is complementary to rather than substitutable for QuantGAN/RSQGAN.

Despite the improvements brought by WGANs, score-based Diffusion models have been increasingly applied to multi-asset return synthesis with the explicit goal of preserving cross-sectional correlation matrices. The empirical case for this claim is, however, mostly \emph{within-paper} rather than cross-paper: CoFinDiff \citep{tanaka2025cofindiff}, Fin-DDPM \citep{takahashi2024generation}, and GBM-Diff \citep{kim2025gbm} each report better cross-correlation preservation against their own GAN baselines on their own chosen datasets, but a shared benchmark that pits all three against the same WGAN/QuantGAN family on the same multi-asset universe is not currently standard practice. The forward diffusion process systematically corrupts real financial data $\mathbf{x}_0$ over $T$ steps by adding Gaussian noise:
\begin{equation}
q(\mathbf{x}_t | \mathbf{x}_{t-1}) = \mathcal{N}(\mathbf{x}_t; \sqrt{1 - \beta_t}\mathbf{x}_{t-1}, \beta_t \mathbf{I})
\end{equation}
where $\beta_t$ dictates the variance schedule. A parameterized neural network $\theta$ then learns the reverse denoising process:
\begin{equation}
p_\theta(\mathbf{x}_{t-1} | \mathbf{x}_t) = \mathcal{N}(\mathbf{x}_{t-1}; \boldsymbol{\mu}_\theta(\mathbf{x}_t, t), \Sigma_\theta(\mathbf{x}_t, t))
\end{equation}
This rigorous score-matching technique allows models like Fin-DDPM and CoFinDiff \citep{takahashi2024generation, tanaka2025cofindiff} to achieve stronger distributional stability than standard GAN-based generators while reducing the risk of mode collapse. To ensure that the generated noise mathematically respects financial reality, models such as GBM-Diff and Fin-Denoiser \citep{kim2025gbm, wang2024denoiser} structurally incorporate Geometric Brownian Motion (GBM) priors directly into the diffusion stochastic differential equations (SDEs), guaranteeing financially sound baseline sequences.

\subsection{Macro-Level Market Simulation: Microstructure and Order Flows}

Macro-level simulation elevates the generative objective from simple price paths to the reproduction of complex market microstructures and multi-agent interactive environments. By generating realistic Limit Order Books (LOB) and order flows, these models facilitate the high-fidelity validation of execution strategies without exposing institutions to real capital risk.

Macro-level market simulation has progressed through three architectural generations, and a careful reading of the empirical evidence suggests that the strongest claims about the diffusion/LMM generation need to be qualified.

The first generation, conditional GANs for correlated price paths, was aimed at reproducing the cross-asset correlation structure that vanilla WGANs miss. \citep{CoMeTS-GAN} explicitly model correlated stock-market time series by injecting a learnable correlation block into the generator, evaluated on the LOBSTER limit-order-book dataset. \citep{istiaque2024cts} extend this idea to \emph{conditional} simulation, conditioning on macro state variables so that the generated regime matches a user-supplied label. \citep{xia2024market} (Market-GAN) go further with a two-stage pipeline in which the first stage generates a market state and the second stage generates price paths conditioned on it. These three papers share a common limitation: their empirical evaluation is mostly on stylized-fact reproduction, and the downstream utility (i.e.\ ``does the synthetic data actually improve trading agents trained on it?'') is rarely reported. \citep{GAN-MoEx} reports a momentum-effect-specific benchmark on Moscow Stock Exchange data and shows GAN-synthetic trajectories reproduce the observed momentum premium; this is one of the few macro-simulation papers with an economic downstream task as the headline metric.

The second generation emerged from macroeconomic conditioning. \citep{alex2024macroeconomic} (MacroSynth) directly condition on macroeconomic factor panels (CPI, rates, credit spreads) so that generated market scenarios are internally consistent with a specified macro state. The reported evaluation is downstream mean-variance and risk-parity allocations computed on synthetic samples, with the synthetic allocations showing sign-consistent risk premia relative to history. The methodological advance here is that the conditioning signal is itself an economic primitive, not a latent regime inferred from the data---a useful design choice for stress-testing use cases where the user specifies the macro scenario in advance.

The third generation, diffusion and large market models, is the most recent wave of macro-simulation work and has moved to diffusion and discrete-token foundation models. \citep{berti2025trades} propose TRADES, a transformer-based denoising diffusion engine conditioned on market state for LOB order-flow simulation. The headline empirical claim is a 3.27- and 3.48-point improvement on the predictive score (MAE between real-data-trained and synthetic-data-trained price-prediction models) over the prior state of the art on two stocks. The authors themselves note that ``there is a notable absence of quantitative metrics for evaluating generative market simulation models in the literature''---the field has been inventing its own evaluation protocols, which makes between-paper comparison fragile. \citep{huang2024diga} (DiGA / DigMA) split the problem into two stages: a conditional diffusion model over time-evolving distribution parameters of mid-price return rate and order arrival rate, and a meta-agent with financial economic priors that samples orders from the learned distributions. This decomposition sidesteps the difficulty of directly diffusing over discrete order sequences. \citep{li2025mars} (MarS) push the boundary furthest by training a Large Market Model (LMM)---a transformer foundation model over discrete order tokens, analogous to a language model over text---and report order-level interactive simulation that is responsive to user-injected orders.

Two critical caveats apply to the Generation-3 claims. First, all three diffusion/LMM papers report improvements over a small set of in-paper baselines, not over a shared benchmark, so the apparent dominance of diffusion over GANs in macro simulation is a within-paper claim rather than a cross-paper empirical fact. Second, the diffusion and LMM papers report evaluation on at most two stocks (TRADES) or single datasets (DiGA, MarS); whether the reported fidelity survives a regime change (e.g.\ 2008 crisis, 2020 COVID) is not established. The survey's earlier claim that ``macro-level LOB simulation and controllable text-to-market generation are still in their infancy'' is therefore better supported than the implicit ``diffusion vastly outperforms GANs'' reading: the field has demonstrably advanced methodologically (transformer denoising, LMM-style discrete tokenization, meta-agent decomposition), but the empirical case for production deployment remains open. We flag this as a measurement gap the field would benefit from closing.

\subsection{Privacy-Preserving and Regulation-Aware Generation}

The financial industry is heavily bottlenecked by strict data-sharing regulations (e.g., GDPR, CCPA). Synthetic data generation provides a mathematically viable pathway to share proprietary datasets across institutional boundaries without violating compliance standards. Models such as NVF-DPGAN \citep{NVF-DPGAN} directly incorporate Differential Privacy (DP) mechanisms into the generative training process. By clipping gradients and injecting calibrated noise during the discriminator's optimization, the framework mathematically guarantees that the generated synthetic sequences cannot be reverse-engineered to leak sensitive information regarding individual retail trading accounts or institutional proprietary alpha signals. The ongoing central research focus in this sub-domain is optimizing the delicate trade-off between strict privacy bounds ($\epsilon$) and the downstream predictive utility of the resulting synthetic series.

% \subsection{Advanced Conditional Generation and Emerging Trends}
\subsection{Emerging Trends}

Recent literature reflects a rapid maturation of financial synthesis, particularly transitioning away from unconditional generation toward highly controllable, context-aware frameworks. Rather than treating synthesis as a purely statistical exercise, recent models simultaneously condition on diverse metadata and partially observed signals to dictate the semantic direction of the synthetic output \citep{gholamrezaei2023conditioning}. 

For instance, \citep{lozano2023dual} propose advanced dual-sourced diffusion networks that align text-based market sentiment directly with numerical trajectory generation, ensuring exceptionally high temporal coherence across modalities. Parallel to improvements in fidelity, there is a critical push to reduce the computational burden of generative sampling. \citep{tao2024} introduce novel distillation techniques specifically designed to accelerate the generation process, a critical requirement for deploying synthetic simulators in live, high-frequency trading environments. Finally, \citep{oppel2025} demonstrate how these advanced conditional generators are practically utilized, embedding them into dynamic stress-testing pipelines to evaluate institutional portfolio resilience under highly specific, synthetically engineered counterfactual crises.

\subsection{Summarization}



Data synthesis currently represents the most dynamic and rapidly evolving branch of Financial Time-Series Generation (FTSG). As summarized in Table \ref{tab:synthesis_comparison}, the field relies on a diverse array of generative architectures, each serving distinct institutional needs. While micro-level return generation has reached a high level of mathematical maturity—reliably capturing volatility clustering and cross-asset correlations via diffusion models—macro-level LOB simulation and controllable text-to-market generation are still in their infancy. Moving forward, the transition towards computationally efficient, privacy-preserving, and condition-specific market simulators will significantly democratize algorithmic trading research and enhance the capability of financial institutions to conduct rigorous, counterfactual stress testing.